% TOP_PROBLEM: {"id": 30006990, "title": "Lex-plus-powers Betti bounds for Artinian quotients over C", "category_id": 4, "category": {"id": 4, "name": "algebra", "display_name": "Algebra", "description": "Group theory, ring theory, field theory, and algebraic structures.", "slug": "algebra", "order_index": 4, "created_at": "2026-07-31T15:26:25.671Z"}, "status": "open"}
% FIRST_POSTED: 2026-10-06
% ENTRY_KIND: statement_only
% !TeX program = lualatex
\documentclass[11pt]{article}
\usepackage[margin=25mm]{geometry}
\usepackage{fontspec}
\setmainfont{Latin Modern Roman}
\usepackage{amsmath,amssymb,mathtools}
\usepackage{unicode-math}
\setmathfont{Latin Modern Math}
\usepackage{xurl,hyperref}
\hypersetup{colorlinks=true,urlcolor=blue}
\setcounter{secnumdepth}{0}
\setlength{\parindent}{0pt}
\setlength{\parskip}{5pt}
\setlength{\emergencystretch}{3em}
\newcommand{\R}{\mathbb R}
\newcommand{\N}{\mathbb N}
\newcommand{\Z}{\mathbb Z}
\newcommand{\Q}{\mathbb Q}
\newcommand{\C}{\mathbb C}
\newcommand{\D}{\mathbb D}
\newcommand{\F}{\mathbb F}
\newcommand{\sref}[2]{\hyperlink{source:#1:#2}{[S#2]}}
\newcommand{\cataloguescope}[1]{#1}
\begin{document}
\section*{Lex-plus-powers Betti bounds for Artinian quotients over C}\label{30006990}
\subsection{Definitions and mathematical statement}
\paragraph{Statement recorded in the supplied source.}
Let $S=C[x_{1},\dots ,x_{n}]$, $n\geq 1$, and let I contain a homogeneous regular sequence of degrees $2\leq a_{1}\leq \cdots \leq a_{n}$. Is there a $lex-plus-(x_{1}^{a_{1}},\dots ,x_{n}^{a_{n}})$ ideal L with the same Hilbert function as I and $\beta ^S_{i,j}(S/I)\leq \beta ^S_{i,j}(S/L)$ for every i,j?
\subsection{Short English statement}
Let $S=C[x_{1},\dots ,x_{n}]$, $n\geq 1$, and let I contain a homogeneous regular sequence of degrees $2\leq a_{1}\leq \cdots \leq a_{n}$. Is there a $lex-plus-(x_{1}^{a_{1}},\dots ,x_{n}^{a_{n}})$ ideal L with the same Hilbert function as I and $\beta ^S_{i,j}(S/I)\leq \beta ^S_{i,j}(S/L)$ for every i,j?
\subsection{Sources}
\begin{itemize}
\item[S1] Original question source or mathematical context.
\url{https://doi.org/10.4171/owr/2009/22}.
\item[S2] OpenAI, manuscript citation page. Source of the formulation motivating this local problem.
\url{https://github.com/openai/math/blob/adc7f1241b42e322a6451854ab7e4b4c146bf78a/preprints/The-Artinian-Lex-Plus-Powers-Betti-Theorem-September-23-2026/README.md}.
\end{itemize}
\end{document}
